\section{总结与延伸}
%% ============================================================

本文系统梳理了 OpenAI Codex 的最佳实践，核心思路可以用一条渐进路径总结：

\begin{importantbox}{Codex 使用的渐进路径}
\textbf{结构化 Prompt} $\rightarrow$ \textbf{AGENTS.md 持久化} $\rightarrow$ \textbf{config.toml 标准化} $\rightarrow$ \textbf{MCP 外部集成} $\rightarrow$ \textbf{Skills 封装} $\rightarrow$ \textbf{Automations 调度}
\end{importantbox}

几个值得深入关注的方向：

\begin{itemize}
    \item \textbf{AGENTS.md 的团队协作模式}：如何在多人团队中维护和演进 \texttt{AGENTS.md}，使其成为活的工程规范文档
    \item \textbf{MCP 生态}：随着 MCP 作为开放标准的发展，更多第三方工具将可通过 MCP 接入 Codex
    \item \textbf{Skill 的 Plugin 化}：从个人本地 skill 到团队共享 plugin 的路径，是提升组织效能的关键
    \item \textbf{Automation 的反思性用法}：不仅用 automation 执行重复任务，还可用于定期回顾和改进工作流本身
\end{itemize}

\subsection{拓展阅读}

\begin{itemize}
    \item OpenAI Codex 官方文档：\url{https://developers.openai.com/codex}
    \item Codex Prompting Guide：\url{https://developers.openai.com/codex/learn}
    \item Model Context Protocol (MCP) 规范：\url{https://modelcontextprotocol.io}
    \item Codex 配置参考：\url{https://developers.openai.com/codex/docs/configuration}
\end{itemize}

%% --- Fin del contenido --- %%

\end{document}
